% ============================================================
% SECTION 7: RESULTS
% ============================================================
% Note: Detailed draft — remove subsections not needed.
% N=23 included sessions unless otherwise stated.
% ============================================================

\section{Results}
\label{sec:results}

Data were collected from 23 participants across two within-subject conditions
(Tempo and Weight), yielding 23 sessions in total. One session was removed prior
to analysis due to a critical equipment fault (participant 04, session 91;
flagged \texttt{DO\_NOT\_USE}), leaving \(N = 23\) sessions in the final
dataset.\footnote{The weight condition has \(N = 22\) because one session
(participant 30) could not be scored for weight owing to insufficient data
coverage.}
Participants were assigned to one of two counterbalanced condition orders
(A-first: Tempo then Weight; B-first: Weight then Tempo).
Descriptive statistics for all outcome measures are reported in
Table~\ref{tab:descriptive}.

% ── Descriptive summary table ────────────────────────────────
\begin{table}[ht]
\centering
\caption{Descriptive statistics for all primary outcome measures
         ($N = 23$ unless noted).}
\label{tab:descriptive}
\begin{tabular}{lcccc}
\hline
\textbf{Measure} & \textbf{$N$} & \textbf{Mean} & \textbf{SD}
                 & \textbf{Range} \\
\hline
\multicolumn{5}{l}{\textit{Tempo condition}} \\
\quad Cadence change (\%) & 19 & 8.0 & 5.3 & [0.3, 22.6] \\
\quad Audio BPM change (\%) & 23 & -- & -- & -- \\
\quad Coupling ratio & 19 & 0.85 & 1.51 & [-2.30, 4.62] \\
\quad Entrainment $r$ (Pearson) & 23 & 0.11 & 0.24 & [-0.24, 0.56] \\
\hline
\multicolumn{5}{l}{\textit{Weight condition}} \\
\quad Stance-phase change (\%) & 11 & 12.7 & 13.7 & [0.0, 46.0] \\
\hline
\multicolumn{5}{l}{\textit{Questionnaires}} \\
\quad Agency (0--10) & 23 & 7.30 & 1.71 & [3.3, 9.7] \\
\quad UEQ-S pragmatic ($-3$ to $+3$) & 22 & 1.26 & 0.99 & [-1.0, 3.0] \\
\quad ARI immersion (1--7) & 23 & 4.54 & 1.00 & [3.0, 6.4] \\
\hline
\end{tabular}
\end{table}

% ── Overview of outcome classification ───────────────────────
Qualitative outcomes were assigned to each session by the primary reviewer
using a five-category scheme:
\textit{influenced\_intended} (cadence/weight shifted in the cued direction and
sustained); \textit{influenced\_transient} (shift detected but did not persist
to the end of the condition); \textit{influenced\_opposite} (shift in the
direction counter to the cue); \textit{no\_effect}; and
\textit{unclear}.\footnote{No sessions were classified as \textit{unclear} in
the final reviewed dataset.}
Table~\ref{tab:outcomes} summarises the distribution of outcomes across both
conditions.

\begin{table}[ht]
\centering
\caption{Outcome classification for Tempo ($N=23$) and Weight ($N=22$)
         conditions.}
\label{tab:outcomes}
\begin{tabular}{lcc}
\hline
\textbf{Outcome} & \textbf{Tempo} & \textbf{Weight} \\
\hline
Influenced --- intended (sustained)  & 4  (17\%)  & 4  (18\%) \\
Influenced --- transient             & 15 (65\%)  & 7  (32\%) \\
Influenced --- opposite direction    & 0  (0\%)   & 3  (14\%) \\
No effect                            & 4  (17\%)  & 8  (36\%) \\
\hline
\textbf{Any influence}               & \textbf{19} (83\%) & \textbf{11} (50\%) \\
\hline
\end{tabular}
\end{table}

% ============================================================
\subsection{RQ1: Tempo Drift and Cadence Change}
\label{sec:results:tempo}
% ============================================================

\subsubsection{Overview of Tempo Influence}

Of the 23 sessions analysed for the Tempo condition, 19 (83\%) showed some
degree of cadence change attributable to the drifting audio tempo. The
majority of influenced participants (15/19; 79\%) showed a transient effect:
their walking cadence shifted toward the audio tempo during the condition but
did not maintain that shift to the end of the trial. Four participants (17\%)
showed a sustained, intended influence, and four (17\%) showed no measurable
effect. No participant walked in the opposite direction to the audio cue in the
Tempo condition.

\subsubsection{Effect Magnitude}

Among influenced participants ($n = 19$), the mean absolute cadence change
relative to baseline was \(8.0\%\) (\textit{SD} = 5.3\%,
range = [0.3\%, 22.6\%]). The largest single effect was observed in P03
(speeding-up condition; 22.6\% cadence increase). The distribution was
right-skewed, with the majority of participants clustered below 10\% change and
a small number of outliers above 13\%.

\subsubsection{Tempo Coupling Ratio}

The tempo coupling ratio --- defined as the participant's cadence change
percentage divided by the audio BPM change percentage --- quantifies how
closely cadence tracked the audio drift on a per-unit basis. For influenced
participants, the mean coupling ratio was 0.85 (\textit{SD} = 1.51). A ratio
of 1.0 would indicate perfect lock-step tracking (cadence changes by exactly
the same proportion as the audio). The large standard deviation reflects
considerable inter-individual variability: some participants over-coupled (e.g.,
P03: 4.62), while others under-coupled or coupled in the opposite sign
(e.g., P01: $-1.31$; P30: $-2.30$). Negative coupling ratios occurred where
the audio slowed the music but the participant's cadence drifted upward (or vice
versa), suggesting the cadence change was not directly driven by audio tracking
in those cases.

\subsubsection{Tempo Entrainment}

Pearson correlation between walking cadence and audio BPM (computed on a 1\,Hz
grid across the condition epoch) served as an index of rhythmic entrainment.
Across all 23 participants the mean entrainment $r$ was 0.11
(\textit{SD} = 0.24, range = [$-0.24$, 0.56]). No participant exceeded the
pre-specified entrainment threshold of $|r| > 0.60$: the maximum observed was
$r = 0.56$ (P39). This indicates that, while some participants showed cadence
changes attributable to the audio, sustained rhythmic entrainment in the strict
sense --- where cadence closely tracks the continuous BPM trajectory --- was not
observed in the present sample.

\subsubsection{Onset of Tempo Effects}

For participants classified as influenced, the temporal onset of the cadence
change was coded into four categories. The majority (15/19; 79\%) showed no
clearly demarcated onset: effects either emerged gradually from the start of
the condition without a discrete inflection point, or were so small as to
preclude reliable onset detection. Two participants showed early onset (change
apparent in the first third of the condition), one showed gradual onset across
the middle third, and one showed late onset (change apparent only in the final
third). This pattern is consistent with the predominantly transient nature of
the effects: responses were diffuse rather than locked to a specific moment.

\subsubsection{Speeding Up vs.\ Slowing Down}

Of the 23 participants, 14 were assigned to the speeding-up condition and 10 to
the slowing-down condition (counterbalanced across order). The overall
proportion of influenced participants was comparable across directions
(speeding-up: 12/14, 86\%; slowing-down: 7/10, 70\%), though the sample is
insufficient to draw reliable directional conclusions.

% ============================================================
\subsection{RQ2: Spectral Weight and Biomechanical Change}
\label{sec:results:weight}
% ============================================================

\subsubsection{Overview of Weight Influence}

Of the 22 sessions with reviewable weight data, 11 (50\%) showed a measurable
biomechanical change. Four sessions (18\%) showed a sustained intended
influence (weight shifted in the cued direction and persisted), seven (32\%)
showed a transient influence, three (14\%) showed a change in the direction
opposite to the cue, and eight (36\%) showed no measurable effect.

The overall influence rate for the Weight condition (50\%) was substantially
lower than for the Tempo condition (83\%).

\subsubsection{Effect Magnitude}

Among participants classified as having an intended or transient influence
($n = 11$), the mean absolute stance-phase change was 12.7\%
(\textit{SD} = 13.7\%, range = [0.0\%, 46.0\%]). The distribution was highly
variable: the majority of influenced participants showed modest changes below
15\%, but two participants showed large effects (P04: 46.0\% stance increase;
P37: 32.5\% stance increase). The three opposite-direction cases (P03: 22.3\%
away from cue; P06: 21.5\%; P32: 47.1\%) were of similar or larger magnitude
to the intended-direction cases, suggesting genuine biomechanical response to
the audio in those participants --- but directed counter to the cue.

\subsubsection{Effect Direction: Heavier vs.\ Lighter}

Participants were cued toward either an increased-weight percept (spectral bass
emphasis; $n = 10$) or a decreased-weight percept (spectral treble/air
emphasis; $n = 12$). The two groups showed markedly different response
profiles:

\begin{itemize}
    \item \textbf{Heavier cue} ($n = 10$): 3 intended, 4 transient, 2 opposite,
          1 no effect. Influence rate: 90\%.
    \item \textbf{Lighter cue} ($n = 12$): 1 intended, 3 transient, 1 opposite,
          7 no effect. Influence rate: 42\%.
\end{itemize}

This asymmetry suggests that participants were more responsive to audio
manipulations designed to evoke a heavier percept than a lighter one. The
heavier condition also produced a higher rate of sustained intended influence
(3/10 vs.\ 1/12). However, the heavier condition also produced more
opposite-direction responses (2/10 vs.\ 1/12), indicating higher overall
responsiveness in both directions.

\subsubsection{Onset of Weight Effects}

Among influenced participants (intended and transient; $n = 11$), onset was
coded as: no identifiable onset (7/11; 64\%), early onset --- change apparent in
the first third (3/11; 27\%), and late onset --- change apparent in the final
third (1/11; 9\%). This pattern broadly mirrors the Tempo condition, with most
effects diffuse rather than sharply triggered.

\subsubsection{Opposite-Direction Responses}

Three participants (P03, P06, P32) showed biomechanical changes in the
direction opposite to their cue. P32 (lighter cue) showed a 47\% increase in
stance-phase duration --- the largest single effect in the entire dataset ---
with an early onset. P03 (lighter cue) and P06 (heavier cue) showed 22\% and
22\% changes respectively. These responses may reflect an active
counterbalancing strategy, individual differences in audio--motor mapping, or
incidental gait variability. Their inclusion in the \textit{influenced\_opposite}
category distinguishes them from the \textit{no\_effect} group on the basis of
change magnitude but treats them as distinct from intended influence for
analysis purposes.

% ============================================================
\subsection{Consciousness Checks and Demand Characteristics}
\label{sec:results:consciousness}
% ============================================================

At four time-points during the session (post-calibration; end of condition 1;
end of condition 2; post-session), participants were asked whether they had
noticed anything different or unusual about the audio. The post-session check
additionally asked them to guess the direction of any perceived tempo or weight
change.

The vast majority of participants reported no awareness of the manipulations.
Across all 23 sessions:

\begin{itemize}
    \item Post-calibration: 1/23 (4\%) reported noticing something.
    \item End of condition 1: 1/23 (4\%) reported noticing something.
    \item End of condition 2: 0/23 (0\%) reported noticing something.
    \item Post-session: 1/23 (4\%) reported noticing something.
\end{itemize}

Of the two participants who provided post-session directional guesses, one
(P32) guessed ``no change'' for both tempo and weight --- despite being
classified as \textit{influenced\_intended} for tempo and
\textit{influenced\_opposite} for weight --- and one (P35) guessed ``slowing
down'' for tempo, which matched their assigned tempo direction. No participant
correctly identified both their tempo and weight cue directions at post-session.
These results suggest that gait changes occurred largely outside participants'
conscious awareness, supporting the ecological validity of the paradigm.

% ============================================================
\subsection{RQ3: Agency, Experience, and Gait Influence}
\label{sec:results:agency}
% ============================================================

\subsubsection{Questionnaire Descriptives}

All 23 participants completed the questionnaire battery. Agency was measured
on a composite 0--10 scale (higher = stronger sense that the audio was
influencing their movement). The mean agency score was 7.30
(\textit{SD} = 1.71, range = [3.3, 9.7]), indicating that most participants
felt at least moderately influenced, and many felt strongly influenced, by the
audio. UEQ-S pragmatic quality scores (scale $-3$ to $+3$) averaged 1.26
(\textit{SD} = 0.99, range = [$-1.0$, 3.0]), consistent with a moderately
positive usability evaluation. ARI immersion scores (scale 1--7) averaged 4.54
(\textit{SD} = 1.00, range = [3.0, 6.4]), indicating moderate-to-good
immersion in the audio experience.

\subsubsection{Correlations Between Subjective Experience and Gait Change}

Spearman rank correlations were computed between each questionnaire measure and
the corresponding objective gait influence magnitude (absolute cadence change
for tempo; absolute stance-phase change for weight). All correlations were
non-significant (Table~\ref{tab:correlations}).

\begin{table}[ht]
\centering
\caption{Spearman rank correlations between subjective measures and objective
         gait influence magnitudes. All $p > .05$.}
\label{tab:correlations}
\begin{tabular}{llcc}
\hline
\textbf{Gait measure} & \textbf{Questionnaire} & \boldmath$r_s$ & \boldmath$p$ \\
\hline
Tempo $|\Delta\%|$  & Agency      & $-0.080$ & .715 \\
Tempo $|\Delta\%|$  & UEQ-S       &  $0.139$ & .527 \\
Tempo $|\Delta\%|$  & ARI         &  $0.076$ & .731 \\
Weight $|\Delta\%|$ & Agency      &  $0.022$ & .924 \\
Weight $|\Delta\%|$ & UEQ-S       & $-0.119$ & .597 \\
Weight $|\Delta\%|$ & ARI         &  $0.212$ & .343 \\
\hline
\end{tabular}
\end{table}

The absence of significant correlations between reported agency and objective
gait change is notable. Participants who reported feeling strongly influenced by
the audio did not exhibit larger measurable gait changes than those reporting
weaker influence, and vice versa. Similarly, immersion (ARI) and pragmatic
usability (UEQ-S) were unrelated to the magnitude of biomechanical response.
This dissociation between subjective experience and objective outcome is
discussed in Section~\ref{sec:discussion}.

% ============================================================
\subsection{Summary of Results}
\label{sec:results:summary}
% ============================================================

The key findings from the three research questions are as follows:

\begin{enumerate}
    \item \textbf{Tempo (RQ1):} Audio tempo drift produced measurable cadence
          changes in 83\% of participants, with mean absolute effect of 8.0\%.
          However, the majority of effects were transient rather than sustained,
          and strict rhythmic entrainment (defined as $|r| > 0.60$) was not
          observed. Cadence--tempo coupling was highly variable across
          individuals (mean ratio = 0.85; \textit{SD} = 1.51).

    \item \textbf{Weight (RQ2):} Spectral weight manipulations produced
          biomechanical changes in 50\% of participants, with mean absolute
          stance-phase change of 12.7\% among those influenced. The heavier cue
          produced more responses (90\%) than the lighter cue (42\%), but also
          more opposite-direction responses. Three participants showed
          counter-directional effects. Onset was predominantly gradual or
          unclassifiable.

    \item \textbf{Agency and Experience (RQ3):} No significant correlations were
          found between any subjective measure (agency, immersion, usability)
          and objective gait influence magnitude. Participants were generally
          unaware of the audio manipulations (0--4\% noticing rates), confirming
          that observed gait changes were not mediated by conscious awareness.
\end{enumerate}
